% Options for packages loaded elsewhere
\PassOptionsToPackage{unicode}{hyperref}
\PassOptionsToPackage{hyphens}{url}
\documentclass[
]{article}
\usepackage{xcolor}
\usepackage[margin=1in]{geometry}
\usepackage{amsmath,amssymb}
\setcounter{secnumdepth}{-\maxdimen} % remove section numbering
\usepackage{iftex}
\ifPDFTeX
  \usepackage[T1]{fontenc}
  \usepackage[utf8]{inputenc}
  \usepackage{textcomp} % provide euro and other symbols
\else % if luatex or xetex
  \usepackage{unicode-math} % this also loads fontspec
  \defaultfontfeatures{Scale=MatchLowercase}
  \defaultfontfeatures[\rmfamily]{Ligatures=TeX,Scale=1}
\fi
\usepackage{lmodern}
\ifPDFTeX\else
  % xetex/luatex font selection
\fi
% Use upquote if available, for straight quotes in verbatim environments
\IfFileExists{upquote.sty}{\usepackage{upquote}}{}
\IfFileExists{microtype.sty}{% use microtype if available
  \usepackage[]{microtype}
  \UseMicrotypeSet[protrusion]{basicmath} % disable protrusion for tt fonts
}{}
\makeatletter
\@ifundefined{KOMAClassName}{% if non-KOMA class
  \IfFileExists{parskip.sty}{%
    \usepackage{parskip}
  }{% else
    \setlength{\parindent}{0pt}
    \setlength{\parskip}{6pt plus 2pt minus 1pt}}
}{% if KOMA class
  \KOMAoptions{parskip=half}}
\makeatother
\usepackage{color}
\usepackage{fancyvrb}
\newcommand{\VerbBar}{|}
\newcommand{\VERB}{\Verb[commandchars=\\\{\}]}
\DefineVerbatimEnvironment{Highlighting}{Verbatim}{commandchars=\\\{\}}
% Add ',fontsize=\small' for more characters per line
\usepackage{framed}
\definecolor{shadecolor}{RGB}{248,248,248}
\newenvironment{Shaded}{\begin{snugshade}}{\end{snugshade}}
\newcommand{\AlertTok}[1]{\textcolor[rgb]{0.94,0.16,0.16}{#1}}
\newcommand{\AnnotationTok}[1]{\textcolor[rgb]{0.56,0.35,0.01}{\textbf{\textit{#1}}}}
\newcommand{\AttributeTok}[1]{\textcolor[rgb]{0.13,0.29,0.53}{#1}}
\newcommand{\BaseNTok}[1]{\textcolor[rgb]{0.00,0.00,0.81}{#1}}
\newcommand{\BuiltInTok}[1]{#1}
\newcommand{\CharTok}[1]{\textcolor[rgb]{0.31,0.60,0.02}{#1}}
\newcommand{\CommentTok}[1]{\textcolor[rgb]{0.56,0.35,0.01}{\textit{#1}}}
\newcommand{\CommentVarTok}[1]{\textcolor[rgb]{0.56,0.35,0.01}{\textbf{\textit{#1}}}}
\newcommand{\ConstantTok}[1]{\textcolor[rgb]{0.56,0.35,0.01}{#1}}
\newcommand{\ControlFlowTok}[1]{\textcolor[rgb]{0.13,0.29,0.53}{\textbf{#1}}}
\newcommand{\DataTypeTok}[1]{\textcolor[rgb]{0.13,0.29,0.53}{#1}}
\newcommand{\DecValTok}[1]{\textcolor[rgb]{0.00,0.00,0.81}{#1}}
\newcommand{\DocumentationTok}[1]{\textcolor[rgb]{0.56,0.35,0.01}{\textbf{\textit{#1}}}}
\newcommand{\ErrorTok}[1]{\textcolor[rgb]{0.64,0.00,0.00}{\textbf{#1}}}
\newcommand{\ExtensionTok}[1]{#1}
\newcommand{\FloatTok}[1]{\textcolor[rgb]{0.00,0.00,0.81}{#1}}
\newcommand{\FunctionTok}[1]{\textcolor[rgb]{0.13,0.29,0.53}{\textbf{#1}}}
\newcommand{\ImportTok}[1]{#1}
\newcommand{\InformationTok}[1]{\textcolor[rgb]{0.56,0.35,0.01}{\textbf{\textit{#1}}}}
\newcommand{\KeywordTok}[1]{\textcolor[rgb]{0.13,0.29,0.53}{\textbf{#1}}}
\newcommand{\NormalTok}[1]{#1}
\newcommand{\OperatorTok}[1]{\textcolor[rgb]{0.81,0.36,0.00}{\textbf{#1}}}
\newcommand{\OtherTok}[1]{\textcolor[rgb]{0.56,0.35,0.01}{#1}}
\newcommand{\PreprocessorTok}[1]{\textcolor[rgb]{0.56,0.35,0.01}{\textit{#1}}}
\newcommand{\RegionMarkerTok}[1]{#1}
\newcommand{\SpecialCharTok}[1]{\textcolor[rgb]{0.81,0.36,0.00}{\textbf{#1}}}
\newcommand{\SpecialStringTok}[1]{\textcolor[rgb]{0.31,0.60,0.02}{#1}}
\newcommand{\StringTok}[1]{\textcolor[rgb]{0.31,0.60,0.02}{#1}}
\newcommand{\VariableTok}[1]{\textcolor[rgb]{0.00,0.00,0.00}{#1}}
\newcommand{\VerbatimStringTok}[1]{\textcolor[rgb]{0.31,0.60,0.02}{#1}}
\newcommand{\WarningTok}[1]{\textcolor[rgb]{0.56,0.35,0.01}{\textbf{\textit{#1}}}}
\usepackage{graphicx}
\makeatletter
\newsavebox\pandoc@box
\newcommand*\pandocbounded[1]{% scales image to fit in text height/width
  \sbox\pandoc@box{#1}%
  \Gscale@div\@tempa{\textheight}{\dimexpr\ht\pandoc@box+\dp\pandoc@box\relax}%
  \Gscale@div\@tempb{\linewidth}{\wd\pandoc@box}%
  \ifdim\@tempb\p@<\@tempa\p@\let\@tempa\@tempb\fi% select the smaller of both
  \ifdim\@tempa\p@<\p@\scalebox{\@tempa}{\usebox\pandoc@box}%
  \else\usebox{\pandoc@box}%
  \fi%
}
% Set default figure placement to htbp
\def\fps@figure{htbp}
\makeatother
\setlength{\emergencystretch}{3em} % prevent overfull lines
\providecommand{\tightlist}{%
  \setlength{\itemsep}{0pt}\setlength{\parskip}{0pt}}
\usepackage{float}
\usepackage{tabularray}
\usepackage[normalem]{ulem}
\usepackage{graphicx}
\usepackage{rotating}
\UseTblrLibrary{booktabs}
\UseTblrLibrary{siunitx}
\NewTableCommand{\tinytableDefineColor}[3]{\definecolor{#1}{#2}{#3}}
\newcommand{\tinytableTabularrayUnderline}[1]{\underline{#1}}
\newcommand{\tinytableTabularrayStrikeout}[1]{\sout{#1}}
\usepackage{bookmark}
\IfFileExists{xurl.sty}{\usepackage{xurl}}{} % add URL line breaks if available
\urlstyle{same}
\hypersetup{
  pdftitle={TidyTuesday Analysis: Chocolate Ratings},
  pdfauthor={Benard\_Adjei},
  hidelinks,
  pdfcreator={LaTeX via pandoc}}

\title{TidyTuesday Analysis: Chocolate Ratings}
\author{Benard\_Adjei}
\date{2026-04-27}

\begin{document}
\maketitle

\section{Why this dataset?}\label{why-this-dataset}

The 2025 January 28 TidyTuesday dataset covers \textbf{water insecurity
across U.S. counties in 2023} specifically the share of households that
lack complete indoor plumbing. At first glance this might look like a
geography or public health dataset, but it's actually a near-perfect
case study for program evaluation. Access to clean water is shaped by
race, poverty, and political geography in ways that are deeply
confounded. You can't just run a regression of ``plumbing access'' on
income and call it a causal estimate there are too many backdoors. This
makes the dataset ideal for applying the tools from this course: DAGs to
map the confounding structure, regression adjustment to close backdoors,
and a discussion of what it would take to actually identify a causal
effect.

\begin{Shaded}
\begin{Highlighting}[]
\CommentTok{\# Load TidyTuesday dataset}
\NormalTok{tuesdata }\OtherTok{\textless{}{-}}\NormalTok{ tidytuesdayR}\SpecialCharTok{::}\FunctionTok{tt\_load}\NormalTok{(}\StringTok{"2025{-}01{-}28"}\NormalTok{)}

\NormalTok{water\_insecurity\_2023 }\OtherTok{\textless{}{-}}\NormalTok{ tuesdata}\SpecialCharTok{$}\NormalTok{water\_insecurity\_2023}
\NormalTok{water }\OtherTok{\textless{}{-}}\NormalTok{ tuesdata[[}\DecValTok{1}\NormalTok{]]}
\end{Highlighting}
\end{Shaded}

\begin{Shaded}
\begin{Highlighting}[]
\NormalTok{water\_clean }\OtherTok{\textless{}{-}}\NormalTok{ water }\SpecialCharTok{|\textgreater{}}
  \FunctionTok{drop\_na}\NormalTok{() }\SpecialCharTok{|\textgreater{}}
  \CommentTok{\# Derive some useful columns for analysis}
  \FunctionTok{mutate}\NormalTok{(}
    \CommentTok{\# Log of total population — large counties span orders of magnitude}
    \AttributeTok{log\_pop =} \FunctionTok{log}\NormalTok{(total\_pop),}
    \CommentTok{\# Bin counties into four population quartiles for grouped comparisons}
    \AttributeTok{pop\_quartile =} \FunctionTok{ntile}\NormalTok{(total\_pop, }\DecValTok{4}\NormalTok{),}
    \AttributeTok{pop\_quartile\_label =} \FunctionTok{factor}\NormalTok{(}
\NormalTok{      pop\_quartile,}
      \AttributeTok{labels =} \FunctionTok{c}\NormalTok{(}\StringTok{"Smallest 25\%"}\NormalTok{, }\StringTok{"25–50\%"}\NormalTok{, }\StringTok{"50–75\%"}\NormalTok{, }\StringTok{"Largest 25\%"}\NormalTok{)}
\NormalTok{    ),}
    \CommentTok{\# Rough Census region from state FIPS — useful for faceting}
    \AttributeTok{state\_fips =} \FunctionTok{as.integer}\NormalTok{(}\FunctionTok{str\_sub}\NormalTok{(geoid, }\DecValTok{1}\NormalTok{, }\DecValTok{2}\NormalTok{)),}
    \AttributeTok{region =} \FunctionTok{case\_when}\NormalTok{(}
\NormalTok{      state\_fips }\SpecialCharTok{\%in\%} \FunctionTok{c}\NormalTok{(}\DecValTok{9}\NormalTok{,}\DecValTok{23}\NormalTok{,}\DecValTok{25}\NormalTok{,}\DecValTok{33}\NormalTok{,}\DecValTok{44}\NormalTok{,}\DecValTok{50}\NormalTok{,}\DecValTok{34}\NormalTok{,}\DecValTok{36}\NormalTok{,}\DecValTok{42}\NormalTok{) }\SpecialCharTok{\textasciitilde{}} \StringTok{"Northeast"}\NormalTok{,}
\NormalTok{      state\_fips }\SpecialCharTok{\%in\%} \FunctionTok{c}\NormalTok{(}\DecValTok{17}\NormalTok{,}\DecValTok{18}\NormalTok{,}\DecValTok{26}\NormalTok{,}\DecValTok{39}\NormalTok{,}\DecValTok{55}\NormalTok{,}\DecValTok{19}\NormalTok{,}\DecValTok{20}\NormalTok{,}\DecValTok{27}\NormalTok{,}\DecValTok{29}\NormalTok{,}\DecValTok{31}\NormalTok{,}\DecValTok{38}\NormalTok{,}\DecValTok{46}\NormalTok{) }\SpecialCharTok{\textasciitilde{}} \StringTok{"Midwest"}\NormalTok{,}
\NormalTok{      state\_fips }\SpecialCharTok{\%in\%} \FunctionTok{c}\NormalTok{(}\DecValTok{10}\NormalTok{,}\DecValTok{11}\NormalTok{,}\DecValTok{12}\NormalTok{,}\DecValTok{13}\NormalTok{,}\DecValTok{24}\NormalTok{,}\DecValTok{37}\NormalTok{,}\DecValTok{45}\NormalTok{,}\DecValTok{51}\NormalTok{,}\DecValTok{54}\NormalTok{,}\DecValTok{1}\NormalTok{,}\DecValTok{21}\NormalTok{,}\DecValTok{28}\NormalTok{,}\DecValTok{47}\NormalTok{,}\DecValTok{5}\NormalTok{,}\DecValTok{22}\NormalTok{,}\DecValTok{40}\NormalTok{,}\DecValTok{48}\NormalTok{) }\SpecialCharTok{\textasciitilde{}} \StringTok{"South"}\NormalTok{,}
\NormalTok{      state\_fips }\SpecialCharTok{\%in\%} \FunctionTok{c}\NormalTok{(}\DecValTok{4}\NormalTok{,}\DecValTok{8}\NormalTok{,}\DecValTok{16}\NormalTok{,}\DecValTok{30}\NormalTok{,}\DecValTok{32}\NormalTok{,}\DecValTok{35}\NormalTok{,}\DecValTok{49}\NormalTok{,}\DecValTok{56}\NormalTok{,}\DecValTok{2}\NormalTok{,}\DecValTok{6}\NormalTok{,}\DecValTok{15}\NormalTok{,}\DecValTok{41}\NormalTok{,}\DecValTok{53}\NormalTok{) }\SpecialCharTok{\textasciitilde{}} \StringTok{"West"}\NormalTok{,}
      \ConstantTok{TRUE} \SpecialCharTok{\textasciitilde{}} \StringTok{"Other"}
\NormalTok{    ) }\SpecialCharTok{|\textgreater{}} \FunctionTok{factor}\NormalTok{(}\AttributeTok{levels =} \FunctionTok{c}\NormalTok{(}\StringTok{"Northeast"}\NormalTok{, }\StringTok{"Midwest"}\NormalTok{, }\StringTok{"South"}\NormalTok{, }\StringTok{"West"}\NormalTok{))}
\NormalTok{  ) }\SpecialCharTok{|\textgreater{}}
  \FunctionTok{filter}\NormalTok{(region }\SpecialCharTok{!=} \StringTok{"Other"}\NormalTok{)   }\CommentTok{\# drops territories}

\FunctionTok{glimpse}\NormalTok{(water\_clean)}
\end{Highlighting}
\end{Shaded}

\begin{verbatim}
## Rows: 835
## Columns: 12
## $ geoid                    <chr> "01069", "04001", "06037", "06097", "06001", ~
## $ name                     <chr> "Houston County, Alabama", "Apache County, Ar~
## $ year                     <dbl> 2022, 2022, 2022, 2022, 2022, 2022, 2022, 202~
## $ geometry                 <chr> "list(list(c(975267.980555021, 975512.9445474~
## $ total_pop                <dbl> 108079, 65432, 9721138, 482650, 1628997, 8978~
## $ plumbing                 <dbl> 93, 2440, 6195, 148, 808, 18, 128, 0, 123, 13~
## $ percent_lacking_plumbing <dbl> 0.08604817, 3.72906223, 0.06372711, 0.0306640~
## $ log_pop                  <dbl> 11.59062, 11.08877, 16.08981, 13.08705, 14.30~
## $ pop_quartile             <int> 2, 1, 4, 4, 4, 1, 2, 2, 3, 4, 4, 3, 4, 3, 1, ~
## $ pop_quartile_label       <fct> 25–50%, Smallest 25%, Largest 25%, Largest 25~
## $ state_fips               <int> 1, 4, 6, 6, 6, 6, 8, 6, 10, 12, 12, 12, 17, 1~
## $ region                   <fct> South, West, West, West, West, West, West, We~
\end{verbatim}

\section{What does the distribution look
like?}\label{what-does-the-distribution-look-like}

Before modeling anything, it's worth just looking at the data.

\begin{Shaded}
\begin{Highlighting}[]
\FunctionTok{ggplot}\NormalTok{(water\_clean, }\FunctionTok{aes}\NormalTok{(}\AttributeTok{x =}\NormalTok{ percent\_lacking\_plumbing, }\AttributeTok{fill =}\NormalTok{ region)) }\SpecialCharTok{+}
  \FunctionTok{geom\_histogram}\NormalTok{(}\AttributeTok{bins =} \DecValTok{60}\NormalTok{, }\AttributeTok{color =} \StringTok{"white"}\NormalTok{, }\AttributeTok{alpha =} \FloatTok{0.85}\NormalTok{) }\SpecialCharTok{+}
  \FunctionTok{facet\_wrap}\NormalTok{(}\SpecialCharTok{\textasciitilde{}}\NormalTok{region, }\AttributeTok{nrow =} \DecValTok{2}\NormalTok{) }\SpecialCharTok{+}
  \FunctionTok{scale\_fill\_viridis\_d}\NormalTok{(}\AttributeTok{option =} \StringTok{"D"}\NormalTok{, }\AttributeTok{end =} \FloatTok{0.85}\NormalTok{) }\SpecialCharTok{+}
  \FunctionTok{scale\_x\_continuous}\NormalTok{(}\AttributeTok{labels =} \FunctionTok{label\_percent}\NormalTok{(}\AttributeTok{scale =} \DecValTok{1}\NormalTok{)) }\SpecialCharTok{+}
  \FunctionTok{labs}\NormalTok{(}
    \AttributeTok{title    =} \StringTok{"Household water insecurity is concentrated in a small fraction of counties"}\NormalTok{,}
    \AttributeTok{subtitle =} \StringTok{"Share of households lacking complete indoor plumbing 2023, by Census region"}\NormalTok{,}
    \AttributeTok{x        =} \StringTok{"\% households lacking plumbing"}\NormalTok{,}
    \AttributeTok{y        =} \StringTok{"Number of counties"}\NormalTok{,}
    \AttributeTok{fill     =} \ConstantTok{NULL}\NormalTok{,}
    \AttributeTok{caption  =} \StringTok{"Source: TidyTuesday 2025{-}01{-}28 | PMAP 8521"}
\NormalTok{  ) }\SpecialCharTok{+}
  \FunctionTok{theme}\NormalTok{(}\AttributeTok{legend.position =} \StringTok{"none"}\NormalTok{)}
\end{Highlighting}
\end{Shaded}

\begin{figure}
\centering
\pandocbounded{\includegraphics[keepaspectratio]{Tidy_Tuesday_files/figure-latex/fig-distribution-1.pdf}}
\caption{Distribution of the share of households lacking complete
plumbing, by Census region. The right skew is severe --- most counties
cluster near zero, but a meaningful tail of counties (especially in the
South and West) have strikingly high rates.}
\end{figure}

\begin{Shaded}
\begin{Highlighting}[]
\NormalTok{water\_clean }\SpecialCharTok{|\textgreater{}}
  \FunctionTok{group\_by}\NormalTok{(region) }\SpecialCharTok{|\textgreater{}}
  \FunctionTok{summarize}\NormalTok{(}
    \AttributeTok{n\_counties    =} \FunctionTok{n}\NormalTok{(),}
    \AttributeTok{median\_pct    =} \FunctionTok{median}\NormalTok{(percent\_lacking\_plumbing),}
    \AttributeTok{mean\_pct      =} \FunctionTok{mean}\NormalTok{(percent\_lacking\_plumbing),}
    \AttributeTok{pct\_above\_2   =} \FunctionTok{mean}\NormalTok{(percent\_lacking\_plumbing }\SpecialCharTok{\textgreater{}} \DecValTok{2}\NormalTok{) }\SpecialCharTok{*} \DecValTok{100}
\NormalTok{  ) }\SpecialCharTok{|\textgreater{}}
  \FunctionTok{tt}\NormalTok{(}
    \AttributeTok{caption =} \StringTok{"Water insecurity summary by Census region"}\NormalTok{,}
    \AttributeTok{digits  =} \DecValTok{2}
\NormalTok{  ) }\SpecialCharTok{|\textgreater{}}
  \FunctionTok{setNames}\NormalTok{(}\FunctionTok{c}\NormalTok{(}\StringTok{"Region"}\NormalTok{, }\StringTok{"Counties"}\NormalTok{, }\StringTok{"Median \%"}\NormalTok{, }\StringTok{"Mean \%"}\NormalTok{, }\StringTok{"\% Counties \textgreater{} 2\%"}\NormalTok{))}
\end{Highlighting}
\end{Shaded}

\begin{table}
\centering
\begin{talltblr}[         %% tabularray outer open
caption={Water insecurity summary by Census region},
]                     %% tabularray outer close
{                     %% tabularray inner open
colspec={Q[]Q[]Q[]Q[]Q[]},
}                     %% tabularray inner close
\toprule
Region & Counties & Median % & Mean % & % Counties > 2% \\ \midrule %% TinyTableHeader
Northeast & 136 & 0.058 & 0.088 & 0 \\
Midwest & 203 & 0.05 & 0.07 & 0 \\
South & 358 & 0.053 & 0.084 & 0 \\
West & 138 & 0.065 & 0.169 & 1.4 \\
\bottomrule
\end{talltblr}
\end{table}

The South and West have the highest average rates and the longest tails.
The median is close to zero everywhere, which means the mean is heavily
pulled by a small group of counties with serious infrastructure
problems.

\section{A DAG for the causal
structure}\label{a-dag-for-the-causal-structure}

This is where the program evaluation framing kicks in. Suppose we want
to ask: \textbf{does population size a proxy for urban infrastructure
investment causally reduce water insecurity?} Or more precisely, what
would happen to a county's plumbing access rate if its population
changed, holding everything else constant?

We can't answer that with a raw correlation, because larger counties
also tend to be wealthier, located in different regions, and have more
political capacity to demand infrastructure investment. All of those
factors independently affect plumbing access. We need to map the
backdoors before we can close them.

\begin{Shaded}
\begin{Highlighting}[]
\NormalTok{water\_dag }\OtherTok{\textless{}{-}} \FunctionTok{dagify}\NormalTok{(}
\NormalTok{  plumbing    }\SpecialCharTok{\textasciitilde{}}\NormalTok{ population }\SpecialCharTok{+}\NormalTok{ region }\SpecialCharTok{+}\NormalTok{ poverty,}
\NormalTok{  population  }\SpecialCharTok{\textasciitilde{}}\NormalTok{ region,}
\NormalTok{  poverty     }\SpecialCharTok{\textasciitilde{}}\NormalTok{ region,}
\NormalTok{  plumbing    }\SpecialCharTok{\textasciitilde{}}\NormalTok{ poverty,}
  \AttributeTok{outcome  =} \StringTok{"plumbing"}\NormalTok{,}
  \AttributeTok{exposure =} \StringTok{"population"}\NormalTok{,}
  \AttributeTok{labels =} \FunctionTok{list}\NormalTok{(}
    \AttributeTok{plumbing   =} \StringTok{"\% Lacking}\SpecialCharTok{\textbackslash{}n}\StringTok{Plumbing"}\NormalTok{,}
    \AttributeTok{population =} \StringTok{"County}\SpecialCharTok{\textbackslash{}n}\StringTok{Population"}\NormalTok{,}
    \AttributeTok{region     =} \StringTok{"Census}\SpecialCharTok{\textbackslash{}n}\StringTok{Region"}\NormalTok{,}
    \AttributeTok{poverty    =} \StringTok{"Economic}\SpecialCharTok{\textbackslash{}n}\StringTok{Conditions"}
\NormalTok{  ),}
  \AttributeTok{coords =} \FunctionTok{list}\NormalTok{(}
    \AttributeTok{x =} \FunctionTok{c}\NormalTok{(}\AttributeTok{population =} \DecValTok{1}\NormalTok{, }\AttributeTok{region =} \FloatTok{2.5}\NormalTok{, }\AttributeTok{poverty =} \FloatTok{2.5}\NormalTok{, }\AttributeTok{plumbing =} \DecValTok{4}\NormalTok{),}
    \AttributeTok{y =} \FunctionTok{c}\NormalTok{(}\AttributeTok{population =} \DecValTok{1}\NormalTok{, }\AttributeTok{region =} \DecValTok{2}\NormalTok{,   }\AttributeTok{poverty =} \DecValTok{0}\NormalTok{,   }\AttributeTok{plumbing =} \DecValTok{1}\NormalTok{)}
\NormalTok{  )}
\NormalTok{)}

\FunctionTok{ggdag\_status}\NormalTok{(water\_dag, }\AttributeTok{use\_labels =} \StringTok{"label"}\NormalTok{, }\AttributeTok{text =} \ConstantTok{FALSE}\NormalTok{) }\SpecialCharTok{+}
  \FunctionTok{theme\_dag}\NormalTok{() }\SpecialCharTok{+}
  \FunctionTok{guides}\NormalTok{(}\AttributeTok{color =} \StringTok{"none"}\NormalTok{) }\SpecialCharTok{+}
  \FunctionTok{labs}\NormalTok{(}
    \AttributeTok{title    =} \StringTok{"Causal structure: population size and water insecurity"}\NormalTok{,}
    \AttributeTok{subtitle =} \StringTok{"Region and economic conditions are backdoor confounders"}
\NormalTok{  )}
\end{Highlighting}
\end{Shaded}

\begin{figure}
\centering
\pandocbounded{\includegraphics[keepaspectratio]{Tidy_Tuesday_files/figure-latex/fig-dag-1.pdf}}
\caption{DAG for the relationship between county population size and
household water insecurity. Region and poverty (proxied here by total
population context) are backdoor confounders they affect both how large
a county is and how much plumbing access it has.}
\end{figure}

\subsection{Interpretating the DAG:}\label{interpretating-the-dag}

The arrow from \textbf{Region → Population} means counties in different
parts of the country tend to differ in size. An arrow from
\textbf{Region → \% Lacking Plumbing} means geography independently
shapes infrastructure. \textbf{Economic Conditions} (poverty) also
connects both to county size and to plumbing access. So a naive
regression of plumbing rates on population size would conflate the
infrastructure investment story with these regional and economic
patterns. To get the causal estimate, we need to condition on region and
economic conditions.

\#The map

A map is the most immediate way to communicate where the problem is
actually located

\begin{Shaded}
\begin{Highlighting}[]
\CommentTok{\# Drop geometry from the sf object so we can join to a pre{-}built county map}
\NormalTok{water\_df }\OtherTok{\textless{}{-}}\NormalTok{ water\_clean }\SpecialCharTok{|\textgreater{}}
  \FunctionTok{st\_drop\_geometry}\NormalTok{() }\SpecialCharTok{|\textgreater{}}
  \FunctionTok{mutate}\NormalTok{(}\AttributeTok{fips =}\NormalTok{ geoid)}

\CommentTok{\# Use map\_data for a lightweight base map (no heavy shapefile needed)}
\NormalTok{counties\_map }\OtherTok{\textless{}{-}} \FunctionTok{map\_data}\NormalTok{(}\StringTok{"county"}\NormalTok{)}

\CommentTok{\# FIPS lookup to join on state + county name}
\FunctionTok{data}\NormalTok{(county.fips, }\AttributeTok{package =} \StringTok{"maps"}\NormalTok{)}
\NormalTok{county\_fips\_clean }\OtherTok{\textless{}{-}}\NormalTok{ county.fips }\SpecialCharTok{|\textgreater{}}
  \FunctionTok{as\_tibble}\NormalTok{() }\SpecialCharTok{|\textgreater{}}
  \FunctionTok{mutate}\NormalTok{(}
    \AttributeTok{fips      =} \FunctionTok{str\_pad}\NormalTok{(}\FunctionTok{as.character}\NormalTok{(fips), }\DecValTok{5}\NormalTok{, }\AttributeTok{pad =} \StringTok{"0"}\NormalTok{),}
    \AttributeTok{polyname  =}\NormalTok{ polyname}
\NormalTok{  )}

\NormalTok{water\_map\_df }\OtherTok{\textless{}{-}}\NormalTok{ water\_df }\SpecialCharTok{|\textgreater{}}
  \FunctionTok{left\_join}\NormalTok{(county\_fips\_clean, }\AttributeTok{by =} \StringTok{"fips"}\NormalTok{) }\SpecialCharTok{|\textgreater{}}
  \FunctionTok{separate}\NormalTok{(polyname, }\AttributeTok{into =} \FunctionTok{c}\NormalTok{(}\StringTok{"region\_map"}\NormalTok{, }\StringTok{"subregion"}\NormalTok{), }\AttributeTok{sep =} \StringTok{","}\NormalTok{, }\AttributeTok{extra =} \StringTok{"drop"}\NormalTok{) }\SpecialCharTok{|\textgreater{}}
  \FunctionTok{right\_join}\NormalTok{(counties\_map, }\AttributeTok{by =} \FunctionTok{c}\NormalTok{(}\StringTok{"region\_map"} \OtherTok{=} \StringTok{"region"}\NormalTok{, }\StringTok{"subregion"} \OtherTok{=} \StringTok{"subregion"}\NormalTok{))}

\FunctionTok{ggplot}\NormalTok{(water\_map\_df, }\FunctionTok{aes}\NormalTok{(}\AttributeTok{x =}\NormalTok{ long, }\AttributeTok{y =}\NormalTok{ lat, }\AttributeTok{group =}\NormalTok{ group,}
                          \AttributeTok{fill =}\NormalTok{ percent\_lacking\_plumbing)) }\SpecialCharTok{+}
  \FunctionTok{geom\_polygon}\NormalTok{(}\AttributeTok{color =} \ConstantTok{NA}\NormalTok{) }\SpecialCharTok{+}
  \FunctionTok{geom\_polygon}\NormalTok{(}
    \AttributeTok{data  =} \FunctionTok{map\_data}\NormalTok{(}\StringTok{"state"}\NormalTok{),}
    \FunctionTok{aes}\NormalTok{(}\AttributeTok{x =}\NormalTok{ long, }\AttributeTok{y =}\NormalTok{ lat, }\AttributeTok{group =}\NormalTok{ group),}
    \AttributeTok{fill  =} \ConstantTok{NA}\NormalTok{, }\AttributeTok{color =} \StringTok{"white"}\NormalTok{, }\AttributeTok{linewidth =} \FloatTok{0.3}\NormalTok{,}
    \AttributeTok{inherit.aes =} \ConstantTok{FALSE}
\NormalTok{  ) }\SpecialCharTok{+}
  \FunctionTok{coord\_map}\NormalTok{(}\StringTok{"albers"}\NormalTok{, }\AttributeTok{lat0 =} \FloatTok{29.5}\NormalTok{, }\AttributeTok{lat1 =} \FloatTok{45.5}\NormalTok{) }\SpecialCharTok{+}
  \FunctionTok{scale\_fill\_viridis\_c}\NormalTok{(}
    \AttributeTok{option   =} \StringTok{"magma"}\NormalTok{,}
    \AttributeTok{direction =} \SpecialCharTok{{-}}\DecValTok{1}\NormalTok{,}
    \AttributeTok{na.value =} \StringTok{"grey88"}\NormalTok{,}
    \AttributeTok{labels   =} \FunctionTok{label\_percent}\NormalTok{(}\AttributeTok{scale =} \DecValTok{1}\NormalTok{),}
    \AttributeTok{breaks   =} \FunctionTok{c}\NormalTok{(}\DecValTok{0}\NormalTok{, }\DecValTok{1}\NormalTok{, }\DecValTok{2}\NormalTok{, }\DecValTok{5}\NormalTok{, }\DecValTok{10}\NormalTok{, }\DecValTok{20}\NormalTok{),}
    \AttributeTok{trans    =} \StringTok{"sqrt"}\NormalTok{,}
    \AttributeTok{name     =} \StringTok{"\% lacking plumbing"}
\NormalTok{  ) }\SpecialCharTok{+}
  \FunctionTok{labs}\NormalTok{(}
    \AttributeTok{title    =} \StringTok{"Where are U.S. households still without complete indoor plumbing?"}\NormalTok{,}
    \AttributeTok{subtitle =} \StringTok{"County{-}level rates, 2023 square root color scale to show variation in low{-}rate counties"}\NormalTok{,}
    \AttributeTok{caption  =} \StringTok{"Source: TidyTuesday 2025{-}01{-}28 | PMAP 8521"}
\NormalTok{  ) }
\end{Highlighting}
\end{Shaded}

\begin{figure}
\centering
\pandocbounded{\includegraphics[keepaspectratio]{Tidy_Tuesday_files/figure-latex/fig-map-1.pdf}}
\caption{County-level household water insecurity across the continental
United States, 2023. Counties shaded darker have higher shares of
households without complete indoor plumbing. Clusters in the rural
South, Appalachia, the Southwest, and parts of the Great Plains are
visible.}
\end{figure}

The geographic clustering is striking. The counties with the highest
rates are not randomly distributed, they're concentrated in specific
places that share histories of underinvestment: the rural South, tribal
lands in the Southwest and Northern Plains, and parts of Appalachia.
This pattern alone tells you there's a regional confounder at work,
exactly as the DAG specified.

\section{Naive vs.~adjusted
regression}\label{naive-vs.-adjusted-regression}

Now for the modeling. Following the DAG, I run two regressions:

\begin{itemize}
\tightlist
\item
  \textbf{Naive}: just regress the plumbing rate on log population,
  ignoring the backdoors.
\item
  \textbf{Adjusted}: add region fixed effects to close the Region
  backdoor path. This is the adjustment-based identification strategy
  from week 7 of the course.
\end{itemize}

\begin{Shaded}
\begin{Highlighting}[]
\CommentTok{\# Naive model — ignores regional confounding}
\NormalTok{model\_naive }\OtherTok{\textless{}{-}} \FunctionTok{lm\_robust}\NormalTok{(}
\NormalTok{  percent\_lacking\_plumbing }\SpecialCharTok{\textasciitilde{}}\NormalTok{ log\_pop,}
  \AttributeTok{data =}\NormalTok{ water\_clean)}
\NormalTok{model\_naive}
\end{Highlighting}
\end{Shaded}

\begin{verbatim}
##                Estimate  Std. Error   t value     Pr(>|t|)    CI Lower
## (Intercept)  0.46014543 0.105539459  4.359937 1.463788e-05  0.25299090
## log_pop     -0.02988451 0.008127945 -3.676761 2.513547e-04 -0.04583817
##                CI Upper  DF
## (Intercept)  0.66729996 833
## log_pop     -0.01393085 833
\end{verbatim}

\begin{Shaded}
\begin{Highlighting}[]
\CommentTok{\# Adjusted model — region fixed effects close the regional backdoor}
\NormalTok{model\_adjusted }\OtherTok{\textless{}{-}} \FunctionTok{lm\_robust}\NormalTok{(}
\NormalTok{  percent\_lacking\_plumbing }\SpecialCharTok{\textasciitilde{}}\NormalTok{ log\_pop }\SpecialCharTok{+}\NormalTok{ region,}
  \AttributeTok{data =}\NormalTok{ water\_clean)}
\NormalTok{model\_adjusted}
\end{Highlighting}
\end{Shaded}

\begin{verbatim}
##                  Estimate Std. Error   t value     Pr(>|t|)     CI Lower
## (Intercept)    0.54136170 0.12789835  4.232750 2.566106e-05  0.290319468
## log_pop       -0.03645625 0.01012895 -3.599213 3.381017e-04 -0.056337626
## regionMidwest -0.03235322 0.01240224 -2.608659 9.253287e-03 -0.056696661
## regionSouth   -0.01543793 0.01146431 -1.346607 1.784744e-01 -0.037940381
## regionWest     0.08174787 0.03810702  2.145218 3.222501e-02  0.006950404
##                   CI Upper  DF
## (Intercept)    0.792403933 830
## log_pop       -0.016574884 830
## regionMidwest -0.008009769 830
## regionSouth    0.007064529 830
## regionWest     0.156545341 830
\end{verbatim}

\begin{Shaded}
\begin{Highlighting}[]
\CommentTok{\# Population{-}quartile model — non{-}parametric version of the same question}
\NormalTok{model\_quartile }\OtherTok{\textless{}{-}} \FunctionTok{lm\_robust}\NormalTok{(}
\NormalTok{  percent\_lacking\_plumbing }\SpecialCharTok{\textasciitilde{}}\NormalTok{ pop\_quartile\_label }\SpecialCharTok{+}\NormalTok{ region,}
  \AttributeTok{data =}\NormalTok{ water\_clean)}
\NormalTok{model\_quartile}
\end{Highlighting}
\end{Shaded}

\begin{verbatim}
##                                  Estimate Std. Error   t value     Pr(>|t|)
## (Intercept)                    0.14131459 0.02383721  5.928321 4.489863e-09
## pop_quartile_label25–50%      -0.03249231 0.02645483 -1.228218 2.197141e-01
## pop_quartile_label50–75%      -0.06118642 0.02491102 -2.456199 1.424587e-02
## pop_quartile_labelLargest 25% -0.08941856 0.02746446 -3.255791 1.176718e-03
## regionMidwest                 -0.02958335 0.01194993 -2.475608 1.349999e-02
## regionSouth                   -0.01408586 0.01116941 -1.261110 2.076244e-01
## regionWest                     0.07992975 0.03748661  2.132221 3.328237e-02
##                                   CI Lower     CI Upper  DF
## (Intercept)                    0.094526137  0.188103051 828
## pop_quartile_label25–50%      -0.084418737  0.019434117 828
## pop_quartile_label50–75%      -0.110082587 -0.012290250 828
## pop_quartile_labelLargest 25% -0.143326715 -0.035510403 828
## regionMidwest                 -0.053039078 -0.006127623 828
## regionSouth                   -0.036009550  0.007837833 828
## regionWest                     0.006349782  0.153509718 828
\end{verbatim}

\begin{Shaded}
\begin{Highlighting}[]
\FunctionTok{modelsummary}\NormalTok{(}
  \FunctionTok{list}\NormalTok{(}
    \StringTok{"Naive"}    \OtherTok{=}\NormalTok{ model\_naive,}
    \StringTok{"Adjusted"} \OtherTok{=}\NormalTok{ model\_adjusted}
\NormalTok{  ),}
  \AttributeTok{stars   =} \ConstantTok{TRUE}\NormalTok{,}
  \AttributeTok{gof\_map =} \FunctionTok{c}\NormalTok{(}\StringTok{"nobs"}\NormalTok{, }\StringTok{"r.squared"}\NormalTok{),}
  \AttributeTok{coef\_rename =} \FunctionTok{c}\NormalTok{(}
    \StringTok{"(Intercept)"}                     \OtherTok{=} \StringTok{"Intercept"}\NormalTok{,}
    \StringTok{"log\_pop"}                         \OtherTok{=} \StringTok{"Log county population"}\NormalTok{,}
    \StringTok{"regionMidwest"}                   \OtherTok{=} \StringTok{"Region: Midwest"}\NormalTok{,}
    \StringTok{"regionSouth"}                     \OtherTok{=} \StringTok{"Region: South"}\NormalTok{,}
    \StringTok{"regionWest"}                      \OtherTok{=} \StringTok{"Region: West"}
\NormalTok{  ),}
  \AttributeTok{title =} \StringTok{"Naive vs. region{-}adjusted estimates: population and water insecurity"}
\NormalTok{)}
\end{Highlighting}
\end{Shaded}

\begin{table}
\centering
\begin{talltblr}[         %% tabularray outer open
caption={Naive vs. region-adjusted estimates: population and water insecurity},
note{}={+ p \num{< 0.1}, * p \num{< 0.05}, ** p \num{< 0.01}, *** p \num{< 0.001}},
]                     %% tabularray outer close
{                     %% tabularray inner open
colspec={Q[]Q[]Q[]},
column{2-3}={}{halign=c,},
column{1}={}{halign=l,},
hline{12}={1-3}{solid, black, 0.05em},
}                     %% tabularray inner close
\toprule
& Naive & Adjusted \\ \midrule %% TinyTableHeader
Intercept & \num{0.460}*** & \num{0.541}*** \\
& (\num{0.106}) & (\num{0.128}) \\
Log county population & \num{-0.030}*** & \num{-0.036}*** \\
& (\num{0.008}) & (\num{0.010}) \\
Region: Midwest &  & \num{-0.032}** \\
&  & (\num{0.012}) \\
Region: South &  & \num{-0.015} \\
&  & (\num{0.011}) \\
Region: West &  & \num{0.082}* \\
&  & (\num{0.038}) \\
Num.Obs. & \num{835} & \num{835} \\
R2 & \num{0.018} & \num{0.053} \\
\bottomrule
\end{talltblr}
\end{table}

The naive model says that larger counties have lower plumbing insecurity
rates which is technically true as a raw correlation. But once we
control for region (closing that backdoor path), the estimate on log
population changes. The regional fixed effects absorb the fact that
larger counties also tend to be in different parts of the country with
different baseline infrastructure. What's left is closer to a within
region comparison: among counties in the same part of the country, do
bigger ones have lower rates?

\section{Visualizing the adjusted
relationship}\label{visualizing-the-adjusted-relationship}

\begin{Shaded}
\begin{Highlighting}[]
\FunctionTok{ggplot}\NormalTok{(water\_clean, }\FunctionTok{aes}\NormalTok{(}\AttributeTok{x =}\NormalTok{ log\_pop, }\AttributeTok{y =}\NormalTok{ percent\_lacking\_plumbing, }\AttributeTok{color =}\NormalTok{ region)) }\SpecialCharTok{+}
  \FunctionTok{geom\_point}\NormalTok{(}\AttributeTok{alpha =} \FloatTok{0.18}\NormalTok{, }\AttributeTok{size =} \FloatTok{0.9}\NormalTok{) }\SpecialCharTok{+}
  \FunctionTok{geom\_smooth}\NormalTok{(}\AttributeTok{method =} \StringTok{"lm"}\NormalTok{, }\AttributeTok{se =} \ConstantTok{TRUE}\NormalTok{, }\AttributeTok{linewidth =} \FloatTok{1.1}\NormalTok{) }\SpecialCharTok{+}
  \FunctionTok{scale\_color\_viridis\_d}\NormalTok{(}\AttributeTok{option =} \StringTok{"D"}\NormalTok{, }\AttributeTok{end =} \FloatTok{0.85}\NormalTok{) }\SpecialCharTok{+}
  \FunctionTok{scale\_y\_continuous}\NormalTok{(}\AttributeTok{labels =} \FunctionTok{label\_percent}\NormalTok{(}\AttributeTok{scale =} \DecValTok{1}\NormalTok{)) }\SpecialCharTok{+}
  \FunctionTok{facet\_wrap}\NormalTok{(}\SpecialCharTok{\textasciitilde{}}\NormalTok{region) }\SpecialCharTok{+}
  \FunctionTok{labs}\NormalTok{(}
    \AttributeTok{title    =} \StringTok{"Larger counties tend to have lower water insecurity but region matters a lot"}\NormalTok{,}
    \AttributeTok{subtitle =} \StringTok{"Each panel shows within{-}region relationship between population size and plumbing access"}\NormalTok{,}
    \AttributeTok{x        =} \StringTok{"Log county population"}\NormalTok{,}
    \AttributeTok{y        =} \StringTok{"\% households lacking plumbing"}\NormalTok{,}
    \AttributeTok{color    =} \ConstantTok{NULL}\NormalTok{,}
    \AttributeTok{caption  =} \StringTok{"Source: TidyTuesday 2025{-}01{-}28 | PMAP 8521"}
\NormalTok{  ) }\SpecialCharTok{+}
  \FunctionTok{theme}\NormalTok{(}\AttributeTok{legend.position =} \StringTok{"none"}\NormalTok{)}
\end{Highlighting}
\end{Shaded}

\begin{verbatim}
## `geom_smooth()` using formula = 'y ~ x'
\end{verbatim}

\pandocbounded{\includegraphics[keepaspectratio]{Tidy_Tuesday_files/figure-latex/unnamed-chunk-1-1.pdf}}

\begin{Shaded}
\begin{Highlighting}[]
\NormalTok{water\_clean }\SpecialCharTok{|\textgreater{}}
  \FunctionTok{group\_by}\NormalTok{(pop\_quartile\_label, region) }\SpecialCharTok{|\textgreater{}}
  \FunctionTok{summarize}\NormalTok{(}
    \AttributeTok{mean\_pct =} \FunctionTok{mean}\NormalTok{(percent\_lacking\_plumbing),}
    \AttributeTok{se\_pct   =} \FunctionTok{sd}\NormalTok{(percent\_lacking\_plumbing) }\SpecialCharTok{/} \FunctionTok{sqrt}\NormalTok{(}\FunctionTok{n}\NormalTok{()),}
    \AttributeTok{n        =} \FunctionTok{n}\NormalTok{()}
\NormalTok{  ) }\SpecialCharTok{|\textgreater{}}
  \FunctionTok{ggplot}\NormalTok{(}\FunctionTok{aes}\NormalTok{(}\AttributeTok{x =}\NormalTok{ pop\_quartile\_label, }\AttributeTok{y =}\NormalTok{ mean\_pct,}
             \AttributeTok{color =}\NormalTok{ region, }\AttributeTok{group =}\NormalTok{ region)) }\SpecialCharTok{+}
  \FunctionTok{geom\_line}\NormalTok{(}\AttributeTok{linewidth =} \DecValTok{1}\NormalTok{) }\SpecialCharTok{+}
  \FunctionTok{geom\_pointrange}\NormalTok{(}\FunctionTok{aes}\NormalTok{(}\AttributeTok{ymin =}\NormalTok{ mean\_pct }\SpecialCharTok{{-}} \FloatTok{1.96} \SpecialCharTok{*}\NormalTok{ se\_pct,}
                       \AttributeTok{ymax =}\NormalTok{ mean\_pct }\SpecialCharTok{+} \FloatTok{1.96} \SpecialCharTok{*}\NormalTok{ se\_pct),}
                  \AttributeTok{size =} \FloatTok{0.5}\NormalTok{) }\SpecialCharTok{+}
  \FunctionTok{scale\_color\_viridis\_d}\NormalTok{(}\AttributeTok{option =} \StringTok{"D"}\NormalTok{, }\AttributeTok{end =} \FloatTok{0.85}\NormalTok{) }\SpecialCharTok{+}
  \FunctionTok{scale\_y\_continuous}\NormalTok{(}\AttributeTok{labels =} \FunctionTok{label\_percent}\NormalTok{(}\AttributeTok{scale =} \DecValTok{1}\NormalTok{)) }\SpecialCharTok{+}
  \FunctionTok{labs}\NormalTok{(}
    \AttributeTok{title    =} \StringTok{"Water insecurity declines with county size but the regional gap persists"}\NormalTok{,}
    \AttributeTok{subtitle =} \StringTok{"Mean \% households lacking plumbing by population quartile and region (±95\% CI)"}\NormalTok{,}
    \AttributeTok{x        =} \StringTok{"County population quartile"}\NormalTok{,}
    \AttributeTok{y        =} \StringTok{"Mean \% lacking plumbing"}\NormalTok{,}
    \AttributeTok{color    =} \StringTok{"Census region"}\NormalTok{,}
    \AttributeTok{caption  =} \StringTok{"Source: TidyTuesday 2025{-}01{-}28 | PMAP 8521"}
\NormalTok{  )}
\end{Highlighting}
\end{Shaded}

\begin{verbatim}
## `summarise()` has grouped output by 'pop_quartile_label'. You can override
## using the `.groups` argument.
\end{verbatim}

\pandocbounded{\includegraphics[keepaspectratio]{Tidy_Tuesday_files/figure-latex/unnamed-chunk-2-1.pdf}}

The quartile plot drives home the confounding story from the DAG. In
every region, the smallest counties tend to have the highest rates of
plumbing insecurity infrastructure investment scales with population.
But the level differs dramatically by region: even the largest Southern
counties have rates comparable to the smallest Northeastern ones. If you
only looked at the raw correlation between population and plumbing
access without accounting for region, you'd get a misleading picture of
what's actually driving the variation.

\section{Putting it together}\label{putting-it-together}

This water insecurity example shows a common issue in program
evaluation: the outcome we care about is influenced by many overlapping
factors. A county's plumbing access isn't just about population size,
it's also shaped by its region, economic conditions, history, and
institutions.

If we run a simple regression without accounting for these factors, we
end up mixing the effect of population with all the other things tied to
it. Controlling for region helps improve the analysis because it removes
one of the most obvious sources of bias and lets us compare counties
more fairly within the same region. It's still not a perfect causal
estimate, but it's clearly better than the naive approach, and using a
DAG helps clarify what we're controlling for and what assumptions we're
making.

If I were taking this further, I'd look for some kind of external
variation in infrastructure investment, like a policy change or funding
rule, that isn't driven by the counties themselves. That's what would
allow us to move from just describing patterns to actually identifying a
causal effect.

\end{document}
